\begin{figure}[H]
\centering
\begin{tikzpicture}[x=2.2cm,y=1.6cm,>=Stealth,font=\small,
    ket/.style={rounded corners=3pt,draw=grisUNI!70,fill=white,
                minimum width=1.25cm,minimum height=0.7cm,inner sep=2pt},
    cero/.style={ket,draw=azulUNI,fill=azulUNI!15,very thick}]

  % rótulos de ejes
  \node[font=\footnotesize] at (0,-1.95) {$m_1\ (\text{partícula }1)$};
  \node[font=\footnotesize,rotate=90] at (-1.75,0) {$m_2\ (\text{partícula }2)$};

  % diagonales M = m1 + m2 = const
  \draw[dashed,grisUNI!70] (0.65,1.35)--(1.35,0.65)   node[right,font=\footnotesize]{$M=2$};
  \draw[dashed,grisUNI!70] (-0.35,1.35)--(1.35,-0.35) node[right,font=\footnotesize]{$M=1$};
  \draw[dashed,azulUNI]    (-1.35,1.35)--(1.35,-1.35) node[right,font=\footnotesize,azulUNI]{$M=0$};
  \draw[dashed,grisUNI!70] (-1.35,0.35)--(0.35,-1.35) node[right,font=\footnotesize]{$M=-1$};
  \draw[dashed,grisUNI!70] (-1.35,-0.65)--(-0.65,-1.35) node[right,font=\footnotesize]{$M=-2$};

  % estados |m1,m2>
  \foreach \a in {-1,0,1}{
    \foreach \b in {-1,0,1}{
      \pgfmathtruncatemacro{\s}{\a+\b}
      \ifnum\s=0
        \node[cero] at (\a,\b) {$|{\a},{\b}\rangle$};
      \else
        \node[ket]  at (\a,\b) {$|{\a},{\b}\rangle$};
      \fi
    }
  }

  % acoplamientos de J^2 dentro de M = 0
  \draw[<->,very thick,rojoUNI] (0.77,-0.77)--(0.23,-0.23)
        node[midway,above right,font=\footnotesize]{$2\hbar^2$};
  \draw[<->,very thick,rojoUNI] (-0.23,0.23)--(-0.77,0.77)
        node[midway,above right,font=\footnotesize]{$2\hbar^2$};

  % matriz del bloque M = 0
  \node[anchor=west,align=left] at (canvas cs:x=4.2cm,y=0.0cm) {%
    $\hat J^{2}\big|_{M=0}=\hbar^2
     \begin{pmatrix}2&2&0\\2&4&2\\0&2&2\end{pmatrix}$\\[8pt]
    \footnotesize autovalores $\{0,\,2\hbar^2,\,6\hbar^2\}$\\
    \footnotesize $\Rightarrow j=0,1,2$\\[6pt]
    \footnotesize $\dim=5+3+1=9$};
\end{tikzpicture}
\caption{Acoplamiento $\ell_1=1\otimes\ell_2=1$ en la base desacoplada
$|m_1,m_2\rangle$. $\hat J^2$ conserva $M=m_1+m_2$, por lo que sólo mezcla
estados sobre la misma diagonal discontinua. En el bloque $M=0$ (azul) los
términos $\hat L_{1\pm}\hat L_{2\mp}$ generan los elementos fuera de la
diagonal $2\hbar^2$ de la ec.~\eqref{eq:matrix_j2_block_zero}.}
\label{fig:cuadricula-m1m2}
\end{figure}
